%% preface.tex — purpose, audience, method, and reading guidance.
\chapter{Preface}

\begin{quotation}
``[T]his arithmetic by 0 and 1 is found to contain the mystery.''

---Gottfried Wilhelm Leibniz, ``Explanation of Binary Arithmetic''
\autocite{leibniz1703binary}
\end{quotation}

Leibniz found mystery in an exact rule. With only zero and one, he could
represent numbers and calculate with them. Three centuries later, the same
two marks describe the words stored in a processor, the instructions that
transform them, and the proofs that say those transformations are correct.
The mystery did not disappear as the rules became precise. Precision gave it
somewhere to live.

This book begins with a modest question: what does it take to establish a
claim about machine code? The question leads farther than its size suggests.
A compiler replaces one instruction sequence with another. A programmer says
that a sequence is the shortest possible. An instruction-set manual prints a
table found by exhaustive search. Each statement concerns a finite machine,
yet words such as \emph{every}, \emph{equivalent}, and \emph{shortest} reach
beyond the examples anyone can inspect.

That reach is one source of the subject's beauty. A short program can still
denote a function on every possible input state in its finite domain. A finite
formula can ask whether any of those inputs makes the program fail. A finite
refutation can sometimes rule out every cheaper program in a declared
language. Small written objects can carry universal claims about machines.
This book studies when they really do.

Testing remains indispensable. It can expose an error in a proposed program.
It cannot by itself show that two programs agree on every input, or that no
cheaper program exists. For those claims we need mathematics, and we need to
know exactly what the machine evidence establishes.

The chapters first build a machine without choosing a commercial instruction
set. They develop words, registers, memory, machine state, operand forms, and
instructions as functions. This foundation lets us compare features often
grouped under the labels reduced instruction set computer (RISC) and complex
instruction set computer (CISC): explicit and implicit operands,
load-store rules and memory operands, fixed and variable encodings, condition
flags, and instruction extensions. RISC-V and x86-64 then become sustained
contrasting companions. Optional vector extensions appear only when the scalar
model's boundary is itself the subject.

After that foundation, the book turns to equivalence, lower bounds, cost
models, proof certificates, and exhaustive search. The deeper subject is the
passage from an executable example to a statement that ranges over every
allowed case. No particular instruction set owns that passage.

Readers should know how programs and processors work. Experience with
a compiler, assembler, or systems code is enough preparation. No prior work
with satisfiability solvers, proof assistants, or proof certificates is
assumed. The mathematics uses functions, finite sets, elementary algebra,
logic, and proof by contradiction. Each piece of notation arrives with the
problem that needs it.

The intended reader is willing to stop at a small example and ask what has
actually been defined. That habit matters more than prior formal training. A
compiler engineer may recognize every instruction in an example yet still
need to state which flags a rewrite preserves. A verification student may
know the solver vocabulary yet still need to follow an effective address into
memory. The book brings those two kinds of care together.

It can be read in three ways. For a first pass, execute the A0 examples by
hand, follow the paired RV64 and x86-64 comparisons, and attempt the exercises
marked by a concrete verb such as Execute, Break, or Transfer. For a course in
semantics, add the definitions and reader proofs, requiring each answer to
name its state and observation. For a course or reading group in
machine-assisted proof, reproduce the artifact routes and their negative
controls. The paths differ in depth, but none should reverse the dependency
order: words before states, states before steps, steps before programs, and
program meaning before optimization.

Examples use small widths because a reader should be able to inspect every
case that carries the explanation. The real architectures retain their real
64-bit widths where the rule depends on them. Moving between those scales is
part of the method. A width-four table reveals a carry boundary; algebra says
which positive widths share it; a bit-vector route checks a declared machine
instance. The text identifies which step supplies which reach.

Figures and colored panels are working parts of the exposition. Definitions
fix objects. Key-idea panels state distinctions worth carrying forward.
Warnings isolate attractive mistakes. Proof kits expose the shape of an
argument before its details. Artifact panels bind a claim to evidence. A
reader should be able to cover any panel and recover its role from the
surrounding prose; the panel organizes the explanation but never replaces it.

Two rules govern the book.

First, scope is part of every technical claim. A statement about an
instruction sequence must name the word width, allowed instructions, operand
forms, constants, and cost model that give the statement its meaning. Change
one of them and the answer may change. Scope is not fine print around a
theorem. It is part of the theorem.

Second, every main theorem has a reader proof and a machine proof. The reader
proof is small enough to reconstruct with pencil and paper. It may use an
invariant, a case split, or a replay. The machine proof runs the complete
stated instance against a stored artifact. One supplies understanding; the
other supplies coverage. A solver output is not an explanation, and an
explanation is not a complete check of a large search space.

The evidence is recorded in objects. Each object joins a statement to its
scope, artifact, checking route, and evidence status. A negative control sends
a nearby false claim or damaged artifact through the same route. The checker
must reject it. This modest demand catches a surprisingly empty ceremony: a
checker that reports success no matter what it receives.

Different objects justify different kinds of belief. Some rely on a solver
verdict. Some carry a certificate that a separate program checks. Others
reproduce a finite calculation or enter a proof-assistant kernel. The text
keeps these routes distinct. A bounded calculation does not become a theorem
for every width, and a solver verdict does not become a kernel proof because
both happen to end successfully.

The exercises are part of the development. They ask the reader to reproduce
a result, damage an artifact, alter a machine, find a boundary, or carry an
argument one step farther. At several points the exercises stop where the
present proof library stops. An honest boundary is not a defect in the
subject. It is where the next piece of mathematics begins.

The accompanying repository contains the object ledger, evidence artifacts,
checkers, and build instructions. The mathematics can be studied without
running the software. Running it reveals another part of the beauty: a claim
about every input eventually becomes particular bytes, read by a particular
program, with a particular way to fail.

The Introduction starts with fixed-width addition and asks which surrounding
state gives that operation its meaning. Chapter 1 then constructs one machine
word. From that smallest exact object, the book builds state, instructions,
memory, programs, two real architectures, and finally the claims we can prove.
